\SetPageTheme{paper}
\pagecolor{white}
\color{black}
\thispagestyle{empty}
\null
\clearpage

\SetPageTheme{paper}
\pagecolor{white}
\color{black}
\thispagestyle{empty}
\begin{center}
{\Huge\bfseries Explicatii finale}\par
\vspace{1em}
{\large Cum poate fi folosit materialul la clasa sau in studiu individual}
\end{center}

\vspace{1.2em}

\begin{tcolorbox}[colback=white,colframe=black!60,boxrule=0.8pt,arc=1.5mm]
\small
\textbf{Pentru profesor}
\begin{itemize}
\item incepeti cu lectura narativa si discutia pe ideea de confidentialitate;
\item insistati pe operatori pana cand efectele lor devin intuitive;
\item nu grabiti trecerea spre algoritm fara trace;
\item folositi exercitiile independente doar dupa stabilizarea schemelor de baza.
\end{itemize}
\end{tcolorbox}

\begin{tcolorbox}[colback=white,colframe=black!60,boxrule=0.8pt,arc=1.5mm]
\small
\textbf{Pentru elev}
\begin{itemize}
\item daca te blochezi, revino la sensul fiecarei variabile;
\item apoi urmareste o singura iteratie a buclei;
\item abia dupa aceea priveste algoritmul ca intreg.
\end{itemize}
\end{tcolorbox}

\clearpage
\pagecolor{cqbg}
\color{cqtext}
\SetPageTheme{challenge}
